\begin{exblock}{复习题 2}
    \item 用记号 $C[a,b]$ 表示在区间 $[a,b]$ 上的连续函数的全体. 设 $f_i \in C[a,b]$, $i = 1, 2, 3$. 定义 $f(x)$ 为三个数 $f_1(x)$, $f_2(x)$, $f_3(x)$ 中介于中间的那一个. 证明: $f \in C[a,b]$.

    \item 设 $f$ 在 $(a,b)$ 上只有第 1 类间断点, 且 $\forall x, y \in (a,b)$ 有不等式
    \[
    f\left(\frac{x+y}{2}\right) \leqslant \frac{f(x) + f(y)}{2}.
\]
    证明: $f \in C(a,b)$.

    \item 对 $n \in \mathbb{N}$, 求满足函数方程
    \[
    f(x + y^n) = f(x) + (f(y))^n,\quad \forall x, y \in \mathbb{R}
\]
    的一切函数 $f$.

    \item 设 $f: \mathbb{R} \to \mathbb{R}$, 且 $f(x^2) = f(x)$, $\forall x \in \mathbb{R}$. 又 $f$ 在 $x = 0$ 与 $x = 1$ 处连续. 证明: $f$ 为常值函数.

    \item 定义函数 $f: [0,1] \to [0,1]$ 如下:
    \[
    f(x) = \begin{cases}
    0.0\, a_1\, 0\, a_2\, 0\, a_3 \cdots, & x = 0.a_1 a_2 a_3 \cdots,\\
    1, & x = 1,
\end{cases}
\]
    其中 $x = 0.a_1 a_2 a_3 \cdots$ 为 10 进制小数表示(当 $x$ 的 10 进制小数表示不惟一时, 一律采用有限小数的方法, 例如 $0.1$ 不表成 $0.0999\cdots$). 试讨论函数 $f$ 的连续性.

    \item 设函数 $f$ 在 $[0,+\infty)$ 上一致连续, 且 $\forall x \in [0,1]$ 有
    \[
    \lim_{n\to+\infty} f(x + n) = 0 \quad (n \in \mathbb{N}).
\]
    证明: $\lim\limits_{x\to+\infty} f(x) = 0$.

    \item 设 $I$ 为区间, 如果存在正的常数 $M$, 使得
    \[
    |f(x) - f(y)| \leqslant M|x - y|,\quad \forall x, y \in I,
\]
    则称 $f$ 在 $I$ 上满足 Lipschitz 条件. 证明: 如果 $f$ 在 $[a,+\infty)$ 满足 Lipschitz 条件 $(a > 0)$, 则 $\frac{f(x)}{x}$ 在 $[a,+\infty)$ 上一致连续.

    \item 证明: 函数 $f$ 在区间 $I$ 上一致连续 $\iff \forall \varepsilon > 0$, $\exists N \in \mathbb{N}$, s.t. 当 $x, y \in I$, $x \neq y$ 且
    \[
    \left|\frac{f(x) - f(y)}{x - y}\right| > N
\]
    时, 恒有 $|f(x) - f(y)| < \varepsilon$.

    \item 证明: 有理函数 $f(x) = \dfrac{1 + x^2}{1 - x^2 + x^4}$ 在 $\mathbb{R}$ 上有界. 进而证明上述有理函数 $f(x)$ 取到最大值 $1 + \dfrac{2}{3}\sqrt{3}$. 而下确界为 $0$, 不可达到.

    \item 设 $f$ 为 $\mathbb{R}$ 上的周期函数, 无最小正周期, 且 $f$ 在某点 $a$ 连续. 证明: $f$ 为常值函数. 换言之, 如果 $f$ 为无最小周期的周期函数且 $f$ 不为常值函数, 则 $f$ 处处不连续. (考察常值函数与 Dirichlet 函数.)

    \item 设函数 $f$ 与 $g$ 在 $[a,b]$ 上连续, 且有 $x_n \in [a,b]$, 使得
    \[
    g(x_n) = f(x_{n+1}),\quad n \in \mathbb{N}.
\]
    证明: 必有一点 $x_0 \in [a,b]$, 使 $f(x_0) = g(x_0)$.

    \item 设函数 $f$ 在区间 $[a,+\infty)$ 上连续有界. 证明: 对每一数 $\lambda$, 存在数列 $x_n \to +\infty$, 使
    \[
    \lim_{n\to+\infty} [f(x_n + \lambda) - f(x_n)] = 0.
\]

    \item 设函数 $f: \mathbb{R} \to \mathbb{R}$ 在有理点上取值为无理数, 在无理点上取值为有理数. 证明: $f$ 不为连续函数.

    \item 设 $f$ 为 $[a,b]$ 上的连续函数, 且在 $\forall [\alpha,\beta] \subset [a,b]$ 上至少有两个不同的最大值点. 证明: $f$ 在 $[a,b]$ 上为常值函数.

    \item 设 $S^2 = \{(x,y,z) \in \mathbb{R}^3 \mid x^2 + y^2 + z^2 = 1\}$ 为单位球面, $f: S^2 \to \mathbb{R}$ 为连续函数. 证明: $\exists \xi \in S^2$, 使 $f(\xi) = f(-\xi)$.

    \item (1) 在半径为 $1$ 的圆内有 $2005$ 个点 $P_1, P_2, \cdots, P_{2005}$. 证明: 在同一平面上一定存在这样一点 $P$, 它到这 $2005$ 个点的距离之和等于 $3000$;

    (2) 在单位球面 $S^2$ 上有 $2005$ 个点 $P_1, P_2, \cdots, P_{2005}$. 证明: 在此球面 $S^2$ 上必有一点 $P_0$, 使
    \[
    \sum_{i=1}^{2005} \widehat{P_0 P_i} = \frac{2005\pi}{2},
\]
    其中 $\widehat{P_0 P_i}$ 表示从 $P_0$ 到 $P_i$ 的球面距离, 即过 $P_0$ 与 $P_i$ 的大圆上, 以 $P_0$ 与 $P_i$ 为端点的劣弧之长.

    \item 设 $f(x)$ 在 $[a,b]$ 上连续, $f(a) < f(b)$. 又设 $\forall x \in (a,b)$, 有
    \[
    \lim_{t\to0} \frac{f(x+t) - f(x-t)}{t} = g(x)
\]
    存在且有限. 证明: $\exists c \in (a,b)$, s.t. $g(c) \geqslant 0$.

    \item 设函数 $f$ 在 $[0,+\infty)$ 上连续, $\forall \alpha > 0$, 有 $\lim\limits_{n\to+\infty} f(n\alpha) = 0$. 证明: $\lim\limits_{x\to+\infty} f(x) = 0$.

    \item 设 $\{f_n\}$ 为 $[a,b]$ 上的连续函数列, $f$ 为 $[a,b]$ 上的连续函数, 它们满足:
    \begin{enumerate}
        \item $f_1(x) \geqslant f_2(x) \geqslant \cdots$, $\forall x \in [a,b]$;
        \item $\lim\limits_{n\to+\infty} f_n(x) = f(x)$, $\forall x \in [a,b]$.
    \end{enumerate}
    证明: 函数序列 $\{f_n\}$ 在 $[a,b]$ 上一致收敛于 $f$.

    \item 设 $a(t)$, $b(t)$ 为 $[0,1]$ 上的连续函数, $0 \leqslant a(t) \leqslant \lambda < 1$. 证明: 方程
    \[
    x = \max_{0\leqslant t\leqslant1} \bigl[b(t) + x a(t)\bigr]
\]
    的解为
    \[
    x = \max_{0\leqslant t\leqslant1} \frac{b(t)}{1 - a(t)}.
\]

    \item (1) 证明: 不存在 $\mathbb{R}$ 上的连续函数, 使它的任一函数值都恰好被取到 $2$ 次;

    (2) 是否存在 $\mathbb{R}$ 上的连续函数, 使它的任一函数值都被恰好取到 $3$ 次.

    \item 设 $f$ 为 $\mathbb{R}$ 上的连续函数, 且 $\lim\limits_{x\to\infty} f \circ f(x) = \infty$. 证明: $\lim\limits_{x\to\infty} f(x) = \infty$.

    \item 应用夹逼定理或 $\ln(1 + x) = x + o(x)\ (x \to 0)$ 证明:
    \[
    \lim_{n\to+\infty} \left(1 + \frac{1^p}{n^{p+1}}\right) \cdots \left(1 + \frac{n^p}{n^{p+1}}\right) = \mathrm{e}^{\frac{1}{p+1}},\quad p > 0.
\]
\end{exblock}
